% BHU PYQ -- Transcribed & Corrected
% Paper: COMMD21: Basic Accounting
% Exam: Under Graduate Semester II Examination, 2025-26 (NEP)
% Department: Faculty of Commerce
% Paper Ref: Conf/even/2025-26/4

\documentclass[11pt,a4paper]{article}
\usepackage[a4paper,top=2.5cm,bottom=2.5cm,left=2.5cm,right=2.5cm,headheight=14pt]{geometry}
\usepackage{amsmath,amssymb}
\usepackage[T1]{fontenc}
\usepackage[utf8]{inputenc}
\usepackage{lmodern,microtype,fancyhdr,enumitem,hyperref,lastpage,parskip,booktabs,tabularx}

\hypersetup{colorlinks=true,urlcolor=black,linkcolor=black,pdftitle={COMMD21: Basic Accounting -- BHU Sem II 2025-26 (NEP MDC)}}
\pagestyle{fancy}\fancyhf{}
\fancyhead[L]{\small\textit{Banaras Hindu University}}
\fancyhead[R]{\small\textit{COMMD21: Basic Accounting (2025-26)}}
\fancyfoot[C]{\small Page \thepage\ of \pageref{LastPage}}

\newlist{parts}{enumerate}{2}
\setlist[parts,1]{label=(\alph*),leftmargin=*,itemsep=4pt,topsep=2pt}
\setlist[parts,2]{label=(\roman*),leftmargin=*,itemsep=3pt,topsep=2pt}

\begin{document}

\begin{center}
  {\large\bfseries BANARAS HINDU UNIVERSITY}\\[2pt]
  {\normalsize Under Graduate Semester II Examination, 2025-26 (NEP)}\\[6pt]
  \rule{0.8\linewidth}{0.5pt}\\[6pt]
  {\large\bfseries FACULTY OF COMMERCE}\\[4pt]
  {\large\bfseries Multidisciplinary Course (MDC)}\\[4pt]
  {\large\bfseries Paper Code: COMMD21 --- Basic Accounting}\\[4pt]
  {\normalsize Time: 2:30 Hours\hfill Maximum Marks: 70}\\[2pt]
  {\small Ref Code: Conf/even/2025-26/4}
\end{center}

\rule{\linewidth}{0.4pt}\smallskip
\noindent\textit{Instruction: Write your Roll No. at the top immediately on the receipt of this question paper.}\\[2pt]
\noindent\textit{Note: Answer all questions. All questions carry equal marks.}
\bigskip

\noindent\textbf{Question 1}
\begin{parts}
  \item Differentiate between Book-keeping and Accounting. Explain the various conventions of Accounting.
\end{parts}

\begin{center}\textbf{OR}\end{center}

Discuss the concepts of Accounting. Give the importance and limitations of Accounting.

\medskip\rule{\linewidth}{0.2pt}\medskip

\noindent\textbf{Question 2}
\begin{parts}
  \item Explain the different kinds of Accounts.
  \item Journalize the following transactions (March 2025):
  \begin{center}
  \small
  \begin{tabular}{r p{9cm} r}
  \toprule
  \textbf{Date (March 2025)} & \textbf{Particulars} & \textbf{Amount (Rs.)} \\
  \midrule
  March 1 & Tomar started business with capital of & 30,000 \\
  March 4 & Purchased goods of & 5,000 \\
  March 7 & Sold goods to Hari & 10,000 \\
  March 11 & Tomar withdrew for personal use & 1,000 \\
  March 15 & Deposited into Bank & 5,000 \\
  March 20 & Paid salary Rs. 1,000 and rent Rs. 500 by issue of cheque & 1,500 \\
  March 25 & Paid to Ratan in full settlement of an account of Rs. 1,000 & 800 \\
  March 31 & Received cash from Tushar (Rs. 1,800) and discount allowed (Rs. 200) & 2,000 \\
  \bottomrule
  \end{tabular}
  \end{center}
\end{parts}

\begin{center}\textbf{OR}\end{center}

Prepare a Cash Book from the following transactions (March 2026):
\begin{center}
\small
\begin{tabular}{r p{9cm} r}
\toprule
\textbf{Date (March 2026)} & \textbf{Particulars} & \textbf{Amount (Rs.)} \\
\midrule
March 1 & Cash in hand Rs. 500 and Cash at Bank & 8,000 \\
March 4 & Received from Bhuvan a cheque of & 2,000 \\
March 7 & Paid insurance premium by cheque & 1,000 \\
March 15 & Purchased goods of Rs. 4,000 in cash @ trade discount of 10\% & 3,600 \\
March 16 & Drew a cheque for office use from bank & 3,000 \\
March 20 & Sold goods in cash & 6,000 \\
March 22 & Settled Account with Mr. Ramesh for Rs. 2,600 by paying cheque of & 2,500 \\
March 24 & Received a cheque on account from Mahesh & 2,000 \\
March 26 & Purchased office stationery in cash & 600 \\
March 27 & Cheque received from Mahesh was returned dishonoured from Bank & 2,000 \\
March 31 & Salary paid by cheque Rs. 3,000. Banked all cash keeping balance of & 1,000 \\
\bottomrule
\end{tabular}
\end{center}

\medskip\rule{\linewidth}{0.2pt}\medskip

\noindent\textbf{Question 3} \\
From the following balances of MH Traders, prepare a Trial Balance for the year ended 31st March 2026:
\begin{center}
\small
\begin{tabular}{l r l r}
\toprule
\textbf{Particulars} & \textbf{Amount (Rs.)} & \textbf{Particulars} & \textbf{Amount (Rs.)} \\
\midrule
Cash in hand & 10,000 & Bank Overdraft & 10,000 \\
Opening Stock & 25,000 & Capital & 80,000 \\
Purchases & 7,500 & Interest Received & 2,000 \\
Sales Return & 5,000 & Sales & 1,61,000 \\
Salary & 15,000 & Purchase Return & 2,500 \\
Debtors & 1,00,000 & Commission Received & 2,500 \\
Machinery & 30,000 & Creditors & 60,000 \\
Wages & 10,000 & Bills Payable & 15,700 \\
Bad Debts & 3,500 & Discount Received & 800 \\
Furniture & 10,000 & General Reserve & 12,000 \\
Carriage Inward & 2,000 & Bills Receivable & 3,500 \\
Office Expenses & 6,000 & Drawings & 12,000 \\
Land \& Building & 22,000 & Rent & 7,500 \\
Insurance & 6,000 & Printing \& Stationery & 4,000 \\
\bottomrule
\end{tabular}
\end{center}

\begin{center}\textbf{OR}\end{center}

What are the various types of Errors in Accounts? Explain with suitable examples. Also explain the process of rectification of errors.

\medskip\rule{\linewidth}{0.2pt}\medskip

\noindent\textbf{Question 4} \\
From the following Trial Balance extracted from the books of N. Kumar, prepare Trading \& Profit \& Loss Account for the year ended 31st Dec. 2025 and prepare Balance Sheet on that date:

\begin{center}
\small
\begin{tabular}{l r l r}
\toprule
\textbf{Debit Particulars} & \textbf{Amount (Rs.)} & \textbf{Credit Particulars} & \textbf{Amount (Rs.)} \\
\midrule
Drawings & 10,000 & Capital & 81,000 \\
Machinery & 60,000 & Creditors & 45,000 \\
Debtors & 40,000 & Sales & 1,40,000 \\
Purchases & 80,000 & Return Outward & 5,000 \\
Return Inward & 4,000 & Discount Received & 500 \\
Wages & 15,000 & Commission & 1,000 \\
Cash in hand & 1,000 & & \\
Cash at Bank & 6,000 & & \\
Salaries & 10,000 & & \\
Repairs & 4,000 & & \\
Rent & 4,500 & & \\
Opening Stock & 20,000 & & \\
Direct Expenses & 5,000 & & \\
Bills Receivable & 10,000 & & \\
Bad Debts & 1,000 & & \\
Carriage Inward & 20,000 & & \\
\midrule
\textbf{Total} & \textbf{2,72,500} & \textbf{Total} & \textbf{2,72,500} \\
\bottomrule
\end{tabular}
\end{center}

\begin{center}\textbf{OR}\end{center}

Give the proforma of Trading and Profit \& Loss Account and Balance Sheet with imaginary figures.

\medskip\rule{\linewidth}{0.2pt}\medskip

\noindent\textbf{Question 5} \\
Write short notes on any four of the following:
\begin{parts}
  \item Users of Accounting Information
  \item Ledger
  \item Suspense Account
  \item Adjustments in Final Accounts
\end{parts}

\end{document}
